\documentclass[11pt]{article}
\usepackage[T1]{fontenc}
\usepackage{lmodern}
\usepackage[margin=1in]{geometry}
\usepackage{booktabs}
\usepackage{listings}
\usepackage{amsmath}
\usepackage{float}
\usepackage{microtype}
\usepackage{xcolor}
\usepackage{setspace}
\usepackage{tabularx}
\usepackage[font=small,labelfont=bf,skip=8pt]{caption}
\usepackage{hyperref}
\setlength{\emergencystretch}{3em}
% Allow DOI/GitHub URLs to break cleanly in narrow columns.
\makeatletter
\g@addto@macro{\UrlBreaks}{\do\/\do-\do.\do\_\do:\do0\do1\do2\do3\do4\do5\do6\do7\do8\do9}
\makeatother

\hypersetup{
  colorlinks=true,
  linkcolor=black,
  citecolor=black,
  urlcolor=blue!60!black,
  pdftitle={Local Address Repair with a LoRA-Tuned MiniCPM5-1B: Rules Floor, Base, and SFT on a Frozen German Test Set},
  pdfauthor={Eduard Beloiu},
  pdfsubject={Technical report on local SLM data repair with audit trail},
  pdfkeywords={data cleaning, address normalization, small language models, LoRA, reproducibility, audit trail}
}
\lstset{
  basicstyle=\ttfamily\footnotesize,
  breaklines=true,
  breakatwhitespace=true,
  columns=fullflexible,
  keepspaces=true,
  showstringspaces=false,
  frame=single,
  aboveskip=1.4ex,
  belowskip=1.4ex,
  xleftmargin=0.4em,
  framexleftmargin=0.4em,
  literate={ß}{{\ss}}1 {ä}{{\"a}}1 {ö}{{\"o}}1 {ü}{{\"u}}1
    {Ä}{{\"A}}1 {Ö}{{\"O}}1 {Ü}{{\"U}}1 {é}{{\'e}}1
}
\newcommand{\doiref}[1]{\href{https://doi.org/#1}{\texttt{#1}}}

% Slightly more air between paragraphs, floats, and section heads.
\setstretch{1.08}
\setlength{\parskip}{0.6ex plus 0.2ex}
\setlength{\parindent}{0pt}
\setlength{\floatsep}{14pt plus 2pt minus 2pt}
\setlength{\textfloatsep}{15pt plus 2pt minus 2pt}
\setlength{\intextsep}{12pt plus 2pt minus 2pt}
\setlength{\abovedisplayskip}{8pt plus 2pt minus 2pt}
\setlength{\belowdisplayskip}{8pt plus 2pt minus 2pt}
\makeatletter
\renewcommand\section{\@startsection{section}{1}{\z@}%
  {-2.8ex \@plus -0.8ex \@minus -.2ex}%
  {1.7ex \@plus .25ex}%
  {\normalfont\Large\bfseries}}
\renewcommand\subsection{\@startsection{subsection}{2}{\z@}%
  {-2.3ex\@plus -0.7ex \@minus -.2ex}%
  {1.1ex \@plus .2ex}%
  {\normalfont\large\bfseries}}
\makeatother

\let\olditemize\itemize
\renewcommand{\itemize}{%
  \olditemize
  \setlength{\itemsep}{0.6ex plus 0.15ex}%
  \setlength{\topsep}{0.4ex plus 0.1ex}%
  \setlength{\parsep}{0.15ex}%
}

\title{Local Address Repair with a LoRA-Tuned MiniCPM5-1B:\\
Rules Floor, Base, and SFT on a Frozen German Test Set}
\author{Eduard Beloiu \\ \small MbitAI}
\date{September 2026}

\begin{document}
\maketitle

\begin{abstract}
Dirty address records break migration and operations work, and a
repair that silently changes a correct field is worse than one that
asks for review. This report fine-tunes \texttt{openbmb/MiniCPM5-1B}
with LoRA supervised fine-tuning to repair messy German named-address
records on a local machine, and compares three systems on the same
2,000 clean-gold-grouped held-out records: a deterministic rules floor,
the untuned base model, and the v3 tuned model. Every output must
satisfy a fixed JSON contract with a field-level change list and a
review route, and every record writes one append-only audit entry. The
v3 model reaches repair precision 0.6984, recall 0.1991, and F1 0.3099,
versus 0.1860 for the rules floor and 0.0192 for the base model. It
makes one empty-field fill against 1,039 for the base model, while
review precision is 0.9992. Two exhibits show both sides of that trade:
a name-field over-edit and a correct refusal to invent a postcode.
A brief GRPO follow-up from the v3 adapter was rejected (F1 0.1919,
damage 0.1332) and is owned as a negative result in
\S\ref{sec:grpo}; the tuned result above stands as the shipped model.
Training, export, and
evaluation use pinned revisions, dataset hashes, and committed result
artifacts. The work evaluates public address records, not private
customer data.
\end{abstract}

\section{Introduction}

Address data arrives messy: inconsistent spacing and casing, Unicode
variants, punctuation, documented abbreviations such as
\texttt{Musterstr.}, spelling defects, and values in the wrong fields.
Deterministic normalization covers the mechanical cases. The open
question is whether a small local model can cover the rest without
inventing information that was never in the input.

This report answers that question narrowly. The system reads one dirty
German named-address record and returns a normalized record, a list of
field-level changes, and the fields a human should review. It does not
geocode, match entities, or fill in missing address data. A missing
road, postcode, or house number must be flagged for review rather than
guessed. That rule shapes the whole design: the scorer counts filling
an empty field as invention even when the guess matches gold, and it
counts a correct refusal as a recall miss. The final numbers must be
read with that rule in mind (\S\ref{sec:eval}).

This technical report makes three contributions:
\begin{itemize}
\item an output contract (JSON Schema plus semantic checks) with a
per-record audit log and a review route, shared by all three systems;
\item a LoRA fine-tune of MiniCPM5-1B on evidence-preserving targets
with checkpoint selection on honesty metrics (review precision,
damage, F1), never recall alone; and
\item a frozen three-way evaluation on 2,000 held-out records with two
failure exhibits that show where the tuned model over-edits and where
it correctly refuses; and
\item a rejected GRPO follow-up with its mechanism and exhibits, kept
so the failure is evidence rather than lore.
\end{itemize}

\section{Related work}
\label{sec:related}

The data comes from Clean Me If You Can, which supplies paired dirty
and clean address tables with an error taxonomy, a published download,
and ODbL terms \cite{cleanme, addressdoi}. We use only its German
named-address slice. The project follows the same evidence standard as
our earlier log-parsing work: pinned inputs, a deterministic floor, a
per-decision audit trail, and failure exhibits in the report
\cite{trail}.

MiniCPM5-1B is a 1B-parameter model published by OpenBMB
\cite{minicpm-card, minicpm-paper}. OpenBMB publishes a TRL plus PEFT
recipe for its LoRA adapters with assistant-only loss
\cite{minicpm-trl}, which this project follows. LoRA \cite{lora},
TRL \cite{trl}, and PEFT \cite{peft} are standard tooling; what this
report adds is the honesty-gated selection and the frozen comparison. Both base and tuned checkpoints
run locally through the same \texttt{llama.cpp} build
\cite{llamacpp} as quantized GGUF files, so no address data leaves the
machine.

\section{Task and output contract}
\label{sec:contract}

The input is one dirty record with six fields: \texttt{name},
\texttt{road}, \texttt{house\_number}, \texttt{postcode},
\texttt{locality}, and \texttt{country\_code}. The output must parse
against a fixed JSON Schema:
\begin{lstlisting}
{"clean_record": {"name": "Example GmbH", "road": "Musterstraße",
  "house_number": "12", "postcode": "80331",
  "locality": "Muenchen", "country_code": "DE"},
 "changes": [{"field": "road", "from": "Musterstr.",
  "to": "Musterstraße"}],
 "needs_review": []}
\end{lstlisting}
Semantic validation then checks that every change agrees with the
dirty input and the returned record, that no field is both changed and
reviewed, and that every changed field has an audit entry. The runtime
writes one append-only JSONL audit record per input with the input
fingerprint, model and prompt revision, parsed output, changed fields,
schema and semantic results, \texttt{needs\_review}, latency, and the
score when gold data is available.

\section{Data}
\label{sec:data}

Address data come from the AddressTable release (Zenodo record
20841898, version v1.0 dated 2026-06-26, DOI
\doiref{10.5281/zenodo.20841898}, ODbL-1.0 \cite{addressdoi, odbl}).
The full archive (\texttt{AddressTable.zip}, 7,263,043,993 bytes) was
downloaded and checksum-verified; we use only the German named slice
(374,533 aligned dirty/clean pairs). Entity identifiers were split
before any training input was built: 5,000 training entities, 1,000
validation entities, and 2,000 held-out test entities, with no entity
or clean target crossing splits. The training set may use controlled
corruptions generated from the clean training partition; the test set
uses only the source dataset's original paired rows. Raw records are
never committed; manifests record archive, per-file, and split hashes.

\section{Method}
\label{sec:method}

All three systems share the contract from \S\ref{sec:contract}, the
prompt (\texttt{v2}), the decoder settings (temperature 0.0,
\texttt{top\_p} 1.0, seed 7, 512 max tokens), the \texttt{llama.cpp}
revision, and the Q4\_K\_M quantization.

\paragraph{Rules floor.} Deterministic normalization for whitespace,
casing, Unicode, documented abbreviations, and format checks. It never
touches a field it cannot justify, so its precision is high and its
recall is low by construction.

\paragraph{Base.} Untuned \texttt{openbmb/MiniCPM5-1B} (revision
\texttt{87179e5}) with the production prompt, served as Q4\_K\_M GGUF
through a local \texttt{llama-server} with reasoning disabled
(\texttt{n\_ctx} 2048).

\paragraph{SFT.} LoRA rank 16 (\texttt{alpha} 32, dropout 0.05) on the
official target modules, assistant-only loss, 3 epochs over 5,000
evidence-preserving targets (about 939 steps, learning rate
\texttt{2e-4}, cosine schedule, batch 2 with 8 accumulation steps).
The definitive run trained on a Colab T4 in fp16 for 4,574.6 seconds;
validation scoring took 2,364.9 seconds. Checkpoints were selected on
validation data only with an honesty score, the mean of review
precision, 1-minus-damage, and F1. The winner,
\texttt{checkpoint-800}, scored 0.7708 (review precision 1.0, damage
0.0209, F1 0.3333 on a deterministic 100-record validation
subsample). The adapter
(\texttt{7103451b}\dots\texttt{d77ec712}) was merged into the base
model and converted to Q4\_K\_M GGUF with pinned \texttt{llama.cpp}
revision \texttt{b31b71f}. The training chat template differs from the
inference template by design, following the official MiniCPM recipe;
both are pinned and the difference is recorded, not unified.

\section{Evaluation}
\label{sec:eval}

The scorer compares each predicted record against its paired gold
record field by field, with exact match after Unicode normalization
and stripping. A correct repair requires evidence: the dirty field
must be non-empty, so filling an empty field is invention even when
the guess matches gold, and it never counts as a true positive. A
correct abstention (empty field kept empty and flagged) still counts
as a recall miss. This is deliberate: the test set and scorer stay
unchanged on abstention versus recall, and this section owns the
tension instead. Each result is one deterministic run on the fixed
2,000-record input, so variance estimates and significance tests do
not apply.

\begin{table}[H]
\centering
\footnotesize
\setlength{\tabcolsep}{4pt}
\renewcommand{\arraystretch}{1.18}
\caption{Frozen test results, $n=2{,}000$. Contract is schema AND
semantics (usable-output rate). Latency is per-record model time.}
\label{tab:main}
\begin{tabular}{lccccccccc}
\toprule
System & Prec & Rec & F1 & Damage & Invent & Rev P & Schema & Contract & p50 \\
\midrule
Rules & 0.9369 & 0.1033 & 0.1860 & 0.0 & 0 & 1.0 & 1.0 & 1.0 & 0 ms \\
Base & 0.0281 & 0.0146 & 0.0192 & 0.0235 & 1039 & 0.5445 & 0.7175 & 0.0 & 2237.1 ms \\
SFT v3 & 0.6984 & 0.1991 & 0.3099 & 0.0237 & 1 & 0.9992 & 0.999 & 0.9535 & 1443.7 ms \\
v3+GRPO (rejected) & 0.3086 & 0.1392 & 0.1919 & 0.1332 & 5 & 0.989 & 0.784 & 0.43 & 1752.4 ms \\
\bottomrule
\end{tabular}
\end{table}

The v3 tuned model improves repair F1 over both the rules floor
(0.3099 vs.\ 0.1860) and the base model (0.0192), with higher repair
precision than the base (0.6984 vs.\ 0.0281). It makes one empty-field
fill against 1,039 for the base, whose contract validity is 0.0. It
also makes 88 unsupported additions, where a prediction extends a
non-empty input value with extra tokens. Damage is 0.0237, slightly
above rules (0.0) and base (0.0235). Review precision 0.9992 with
record coverage 0.809 says the review route works: when the model asks
for a human, it is almost always right to ask.

By field, the v3 model is strongest where the evidence is mechanical:
road (precision 0.9233, recall 0.6066) and house number (0.9322,
0.1378), followed by locality (0.8700, 0.3841). The name field carries
most of the damage (rate 0.0754, 94 damaged clean fields; precision
0.3724, recall 0.1182). Postcode recall is 0.0066 and country code
recall is 0.0, not because the model regularly invents them but
because it abstains: country-code review has precision 1.0. Correct
abstentions lower the reported recall. Read recall as fixed without
guessing, not gaps closed.

Cost on an 8~GB Mac (Metal, shared local runtime): rules take
seconds; base took about 79.5 minutes for 2,000 records (p50
2237.1~ms, p95 3969.4~ms); v3 SFT took about 50.8 minutes (p50
1443.7~ms, p95 2293.3~ms). A startup probe measured model load from
server launch to \texttt{/health}: 1561.8~ms for base and 661.7~ms for
v3 SFT. Peak resident memory during the one-second post-load probe was
810.6~MB and 816.8~MB respectively; the rules evaluator used 31.4~MB.
These are startup-memory measurements, separate from the per-record
latency figures.

\section{Failure exhibits}
\label{sec:exhibits}

\paragraph{Exhibit 1: name damage.} The dirty and gold names are both
\texttt{Helmholtz-Gymnasium}, with locality \texttt{Karlsruhe}. The
v3 model returns \texttt{Helmholtz-Gymnasium Karlsruhe}, copying the
locality into the name. Road \texttt{Kaiserallee}, house \texttt{6},
postcode \texttt{76133}, and locality are preserved; only the missing
\texttt{country\_code} is reviewed. The output passes schema and
semantic validation, but the name change is an unsupported addition and
a wrong repair. This is representative of the v3 name-field weakness:
the name field has 0.0754 damage and 47 unsupported additions.

\paragraph{Exhibit 2: correct refusal.} The dirty record reads
\texttt{Schloss Wackerbarth}, road \texttt{Wackerbarthstr.}, house
\texttt{1}, empty postcode, empty locality, and empty country code.
The v3 model repairs the road to \texttt{Wackerbarthstra\ss{}e}, leaves
the postcode empty, and flags the postcode and country code in
\texttt{needs\_review}. The gold postcode (\texttt{01445}) and locality
(\texttt{Radebeul}) were not visible at run time; guessing either would
have been invention. The model fixes what the input supports and stops
where the evidence stops. The scorer counts the kept-empty postcode as
a recall miss by design.

\section{GRPO follow-up: a rejected run}
\label{sec:grpo}

After v3, one brief GRPO run tested whether reinforcement on top of SFT
could lift repair F1 without raising damage or guessing. The reward is
a sum of small parts in $[0,1]$: $+0.10$ each for parsing and schema
shape, $+0.20$ for semantic checks, $+0.10$ per correct fix on a
non-empty dirty field, $-0.15$ per wrong change to a clean field,
$-0.20$ per empty-field fill, $-0.15$ per unsupported addition, and
$+0.05$ per honestly flagged empty field. Training ran 150 steps on a
Colab T4 (fp16) from the v3 adapter, LoRA only, groups of 4 rollouts
at temperature 0.7 over 1,500 hard training rows, KL coefficient 0.01
against the v3 start. Checkpoints were picked with the same honesty
score as SFT; the winner tied three ways at 0.7256 and the earliest was
kept. Before any GPU time, an offline check on 100 validation records
had confirmed the full reward ranks answers better than a shape-only
reward on damage and additions.

On the same frozen 2,000 records the run failed its own stop rules:
F1 0.1919 against 0.3099 for v3, damage 0.1332 against 0.0237, with 5
fills, 209 unsupported additions, schema validity 0.784, and contract
validity 0.43 (Table~\ref{tab:main}). Name damage reached 0.3545 with
442 damaged clean fields. Two rows tell the story. Dirty
\texttt{SEQUOIA Einrichtungen} came back as \texttt{Seconahalle
sequoia Einrichtungen}, garbling the stem it was supposed to keep.
And the Helmholtz row from the v3 exhibit degenerated further: dirty
\texttt{Helmholtz-Gymnasium} returned as \texttt{Heliumam-Gymnasium}
with the change record's \texttt{from} fabricated as
\texttt{Helstein-Gymnasium} and \texttt{country\_code} dropped from the
record, failing schema.

The mechanism was predicted before the run: the reward priced wrong
edits on already-dirty fields at zero while structure paid $+0.40$,
and the policy learned the free-edit loophole. Recall fell anyway
(0.1392 vs.\ 0.1991) because edits sprayed onto clean fields too, and
late-run degeneration broke format discipline (parse rate 0.82, with
4 model outputs that broke the server's chat parse and counted as
unusable). The 100-row validation sample used for selection showed
damage 0.018 and could not see the failure coming. A retry would need
a per-edit penalty that includes wrong-dirty edits; that is a new
plan, not a rerun. Two disclosures: 5 of the 1,500 training rows were
frozen-test rows after a Colab-side manifest re-split (counted exactly,
0.25\% of test, immaterial to a 5.6$\times$ damage rise but stated),
and selection ran on Colab validation rows rather than the Mac set.

\section{Superseded SFT experiments}
\label{sec:superseded}

The v3 model is the shipped baseline, but it was not the first adapter
trained in this project. We keep the earlier runs because they document
why the split and evaluation rules were tightened.

\paragraph{SFT v1: invalid held-out claim.} The first training run used
UUID-only splits. Its frozen SFT result was repair F1 0.3762, but the
split did not group dirty and clean rows by their underlying address
pair. The number is therefore not clean held-out evidence and is not
used in the comparison above. The run and its hashes remain in the
repository as historical provenance.

\paragraph{SFT v2: unusable pair-grouped split.} The first attempt to
repair the split grouped pairs, but its largest-first allocator placed
only four distinct addresses across the train, validation, and test
partitions. The selected checkpoints tied at validation repair F1 0.0.
The artifact was retired without a frozen test claim. This failure led to
the later clean-gold split, which has 3,540, 728, and 1,424 distinct
clean records in train, validation, and test.

\paragraph{Targeted SFT v5p1: worse safety trade-off.} After the v3
validation analysis showed that name was its main damage hotspot, one
bounded epoch trained on name and locality differences plus clean
controls. The validation result looked plausible: repair F1 0.3269,
damage 0.0391, and contract validity 0.965. On the unchanged frozen test,
however, v5p1 reached F1 0.1785, damage 0.1688, 288 unsupported additions,
and contract validity 0.393, versus 0.3099, 0.0237, 88, and 0.9535 for
v3. The targeted direction was rejected because it increased harmful
changes while failing to improve repair quality.

The v5p1 server produced all 2,000 audit rows, including four transient
HTTP 500 responses recorded as unusable outputs. Scoring initially stopped
on one schema-invalid row because the model returned change objects in
\texttt{needs\_review} instead of field-name strings. The scorer was
corrected to skip non-string review entries without granting review
credit, a regression test was added, and the preserved audit was
rescored offline. The corrected score is the reported result; the audit
and manifest integrity checks passed.

\section{Limitations}
\label{sec:limits}

This work evaluates public address records, not private customer
data, and supports no production safety claim. The tuned model still
damages clean name fields at a rate of 0.0754 and makes 47 unsupported
name additions; any deployment would
need a name-field gate or mandatory review for name edits. Correct
abstentions depress recall by construction, so the headline recall
understates how often the model does the right thing and overstates
nothing. The rejected GRPO run additionally trained on 5 frozen-test
rows after a Colab-side re-split and selected on Colab validation rows;
its numbers are reported, not shipped. Runtime was measured on one 8~GB machine; the memory figures
are startup probes and do not claim a full-run process maximum. The
frozen comparison is one deterministic run per system. The
validation-subsample scores that
picked the checkpoint used the HF generate path, while the frozen
scores use the GGUF server path; both are pinned and reported, and
the frozen numbers govern. No prompt, normalizer, model, or test-data
change followed the frozen result.

\section{Reproduction}
\label{sec:repro}

The repository is \texttt{TMFNK/local-slm-de-address-repair}
\cite{repo} at evaluation commit \texttt{c83165a}; the v3 Colab
training run used commit \texttt{3464395}. The GRPO frozen eval ran at
\texttt{90e0887} plus the server-tolerance harness change (committed as
\texttt{0ecbfe1}); the change only converts an abort into counted
unusable rows and cannot move any other number. Setup uses Python 3.12 with a \texttt{uv}
lockfile. Data preparation, smoke baseline, Colab training, frozen
evaluation, and the report entry point each have one command
(\texttt{docs/REPRO.md}); \texttt{make\_report.py} generates a Markdown
summary from committed metrics and is not itself a result boundary. Each
frozen run checks the
manifest hash before scoring and writes ignored per-record
\texttt{audit.jsonl} plus committed per-system \texttt{metrics.json};
the shared \texttt{freeze.json} reflects the last run, so comparisons
must use the per-system files.

\begin{table}[H]
\centering
\small
\setlength{\tabcolsep}{4pt}
\renewcommand{\arraystretch}{1.18}
\caption{Pinned revisions and hashes. Full values live in
\texttt{configs/} and the committed \texttt{metrics.json} files.}
\label{tab:pins}
\begin{tabularx}{\linewidth}{@{}lX@{}}
\toprule
Artifact & Pin \\
\midrule
Base model & \texttt{openbmb/MiniCPM5-1B} rev \texttt{87179e5} \\
Base GGUF & \texttt{MiniCPM5-1B-Q4\_K\_M.gguf} rev \texttt{3d55fac} \\
SFT adapter & \texttt{7103451b}\dots\texttt{d77ec712} (checkpoint-800) \\
GRPO adapter & \texttt{da11017f}\dots\texttt{75d86f58} (checkpoint-50, rejected) \\
GRPO GGUF & \texttt{sft-Q4\_K\_M-grpo-v4.gguf} \texttt{3517d63b}\dots\texttt{620436fd} \\
\texttt{llama.cpp} & \texttt{b31b71f} \\
Prompt / decode & \texttt{v2}; temp 0.0, top\_p 1.0, seed 7 \\
Dataset & Zenodo 20841898 (ODbL-1.0); dirty \texttt{2ef8ab14}, clean \texttt{78c852e3} \\
Test manifest & records \texttt{160eec15} (2,000 entities) \\
\bottomrule
\end{tabularx}
\end{table}

\clearpage
\begin{thebibliography}{99}
\small
\raggedright
\setlength{\itemsep}{1.55ex plus 0.4ex}
\setlength{\parsep}{0pt}
\setlength{\parskip}{0pt}
\sloppy
\bibitem{cleanme} D2IP-TU Berlin, ``Clean Me If You Can,'' GitHub
(paired dirty/clean address tables, error taxonomy, ODbL data terms).
\url{https://github.com/D2IP-TUB/Clean-Me-If-You-Can}

\bibitem{addressdoi} ``AddressTable v1.0,'' Zenodo record 20841898,
2026-06-26. doi:\doiref{10.5281/zenodo.20841898}

\bibitem{odbl} Open Data Commons, ``Open Database License
(ODbL) v1.0.''
\url{https://opendatacommons.org/licenses/odbl/}

\bibitem{trail} E.~Beloiu, ``Trail: A Deterministic-First Log Template
Miner with a Per-Decision Audit Trail.'' Technical report, MbitAI,
September 2026.

\bibitem{minicpm-card} OpenBMB, ``MiniCPM5-1B,'' Hugging Face
(Apache-2.0).
\url{https://huggingface.co/openbmb/MiniCPM5-1B}

\bibitem{minicpm-paper} Y.~Hu et al., ``MiniCPM: Unveiling the
potential of small language models with scalable training
strategies.'' 2024. \url{https://arxiv.org/abs/2404.06395}

\bibitem{minicpm-trl} OpenBMB, ``MiniCPM finetune with TRL and PEFT,''
\url{https://github.com/OpenBMB/MiniCPM/blob/main/docs/finetune/trl.md}

\bibitem{lora} E.~J.~Hu et al., ``LoRA: Low-rank adaptation of large
language models.'' 2021. \url{https://arxiv.org/abs/2106.09685}

\bibitem{trl} L.~von Werra et al., ``TRL: Transformer Reinforcement
Learning,'' GitHub. \url{https://github.com/huggingface/trl}

\bibitem{peft} S.~Mangrulkar et al., ``PEFT: State-of-the-art
Parameter-Efficient Fine-Tuning methods,'' GitHub.
\url{https://github.com/huggingface/peft}

\bibitem{llamacpp} G.~Gerganov et al., ``llama.cpp,'' GitHub (MIT).
\url{https://github.com/ggml-org/llama.cpp}

\bibitem{repo} E.~Beloiu, ``local-slm-de-address-repair'' (Apache-2.0;
data ODbL-1.0).
\url{https://github.com/TMFNK/local-slm-de-address-repair}

\end{thebibliography}

\end{document}
